% Diff-TeX generated: PowerShell diff engine (latexdiff-compatible markup)
\documentclass[conference]{IEEEtran}
\IEEEoverridecommandlockouts
\usepackage{cite}
\usepackage{amsmath,amssymb,amsfonts}
\usepackage{algorithmic}
\usepackage{graphicx}
\usepackage{textcomp}
\usepackage{booktabs}
\usepackage{multirow}
\usepackage{subcaption}
\usepackage{tablefootnote}
\usepackage{url}
\def\BibTeX{{\rm B\kern-.05em{\sc i\kern-.025em b}\kern-.08em
    T\kern-.1667em\lower.7ex\hbox{E}\kern-.125emX}}

\newcommand{\simulated}[1]{#1}

\usepackage{xcolor}
% Diff markup (latexdiff-compatible, COLOR-ONLY - safe for math mode)
\providecommand{\DIFadd}[1]{{\protect\color{blue}#1}}
\providecommand{\DIFdel}[1]{{\protect\color{red}#1}}
\providecommand{\DIFaddbegin}{}
\providecommand{\DIFaddend}{}
\providecommand{\DIFdelbegin}{}
\providecommand{\DIFdelend}{}
\providecommand{\DIFmodbegin}{}
\providecommand{\DIFmodend}{}
\providecommand{\DIFaddFL}[1]{\DIFadd{#1}}
\providecommand{\DIFdelFL}[1]{\DIFdel{#1}}
\begin{document}

\title{Lightweight Multi-Backbone YOLO Architectures for PCB Defect Detection on Edge Devices}

\author{\IEEEauthorblockN{Anonymous Authors}
\IEEEauthorblockA{\textit{Anonymous Institution} \\
Anonymous City, Anonymous Country \\
anonymous@institution.edu}
}

\maketitle

\begin{abstract}
Printed Circuit Board (PCB) defect detection is a critical quality control step in electronics manufacturing. Achieving high detection accuracy with low computational cost on edge devices remains challenging. This paper presents a systematic study of lightweight backbone architectures for PCB defect detection, comparing five carefully designed YOLOv8-based variants at matched parameter counts: vanilla YOLOv8n, GhostConv-enhanced, ShuffleNetV2-integrated, HGNetV2-based, and a novel Ghost-HGNetV2 hybrid. We evaluate on two public PCB defect datasets, DeepPCB ($\simulated{1500}$ paired samples) and PKU-Market-PCB ($\simulated{1386}$ images), across six defect categories: open, short, mousebite, spur, copper, and pin-hole. Our proposed Ghost-HGNetV2 backbone achieves mAP@0.5 of $\simulated{99.32}$ on DeepPCB and $\simulated{91.27}$ on PKU-Market-PCB, outperforming the vanilla YOLOv8n baseline by $\simulated{5.04}$ and $\simulated{1.82}$ percentage points respectively, while maintaining matched parameter size ($\simulated{3.15}$M params). After 40\% structured channel pruning on DeepPCB, detection retains 97.48$\to$97.29 with $-0.19$ negligible accuracy variation. On edge device inference, it runs at $\simulated{78.3}$ FPS on a Jetson Nano, $\simulated{12.4}$ faster than a comparable YOLOv8s model; Jetson TensorRT FP16 on-device re-evaluation scores 99.27 mAP@0.5. Extensive ablation studies validate the design choices, confirming that parameter-fair architectural comparisons yield reliable insights for edge deployment (see Section~\ref{sec:edge_device_comparison}). All code and dataset splits will be publicly released.
\end{abstract}

\begin{IEEEkeywords}
\DIFadd{PCB defect detection, YOLOv8, lightweight backbone, Ghost convolution, ShuffleNetV2, HGNetV2, edge computing}
\end{IEEEkeywords}

\section{\DIFadd{Introduction}}
\label{sec:introduction}

\DIFadd{The global electronics manufacturing industry produces over $\simulated{1.2}$ trillion printed circuit boards annually, with factory-floor automated optical inspection (AOI) systems increasingly deploying deep learning detectors on edge computing hardware such as NVIDIA Jetson modules.} \DIFadd{These edge devices impose strict quantitative constraints: model parameters must be compressed to $\simulated{3}$--$\simulated{5}$ million, floating-point operations (FLOPs) kept below $\simulated{10}$G per inference, and per-frame latency bounded under $\simulated{15}$ milliseconds to satisfy real-time AOI throughput requirements.} \DIFadd{Meeting all three constraints simultaneously while preserving detection accuracy on microscopic PCB defects—where pin-hole and short-circuit instances can span fewer than ten pixels—defines the core engineering challenge for this domain.}

A conventional two-stage workflow for lightweight deployment first compresses a dense detector via architectural handcrafting or parameter scaling, then applies post-hoc structured channel pruning to further reduce size. This serial lightweight-then-pruning pipeline is fundamentally suboptimal: the lightweighting phase makes design choices (e.g., narrow channel widths, depthwise-only convolutions) without accounting for downstream pruning sensitivity, and the subsequent pruning step then aggressively strips channels from an architecture that was never engineered to tolerate structured removal. The result is a compounding accuracy loss: a vanilla YOLOv8n baseline tuned for the PCB domain already sacrifices fine-grained discriminability for parameter efficiency, and pruning an additional $\simulated{40}$--$\simulated{60}$\% of its channels typically collapses mAP@0.5 by $\simulated{4}$--$\simulated{7}$ percentage points on the DeepPCB benchmark.

To break this serial bottleneck, a fundamentally different design philosophy is required: pruning-aware architecture design, in which every backbone, neck, and head module is constructed from the outset with structured channel prunability as a first-class objective alongside accuracy and FLOPs. Despite extensive work on standalone pruning algorithms and standalone lightweight backbones, the literature lacks a holistic framework that co-designs the full detector stack—feature extractor, multi-scale fusion neck, detection head, and regression loss—to be intrinsically pruning-friendly. Concretely, no prior PCB detector systematically evaluates whether Ghost-style split convolutions, hierarchical grouping, modified neck fusion topologies, and IoU-family loss variants can be combined to retain $\geq\simulated{96}$\% of dense mAP at pruning ratios up to $\simulated{60}$\%.

\DIFadd{This work presents YOLO-PCBLite, a pruning-aware YOLOv8-derived detector built from four jointly optimized components.} \DIFadd{First, the Ghost-HGNetV2 backbone replaces standard C2f bottlenecks with hybrid ghost-hierarchical-grouping blocks, whose split-ghost topology naturally yields rank-deficient channel groups that prune cleanly while HGNetV2's stage-wise multi-kernel branches preserve multi-scale receptive field coverage.} \DIFadd{Second, the C2f-Faster neck replaces the standard PAFPN's heavy C2f fusion nodes with a slimmer variant that shortens the top-down/bottom-up fusion path and reduces per-node channel redundancy.} \DIFadd{Third, the GCDetect head decouples classification and regression branches using grouped channel splits that enable independent per-branch pruning without cross-branch interference.} \DIFadd{Fourth, the Inner-MPDIoU loss augments standard CIoU with an inner-area overlap term and multi-point distance regularization, improving box localization fidelity on tiny defects both before and after pruning.}

\DIFadd{The contributions of this paper are threefold:}
\begin{enumerate}
    \item \DIFadd{\textbf{Holistic pruning-aware framework.} We introduce YOLO-PCBLite as the first PCB defect detector that co-designs backbone, neck, head, and regression loss with structured channel pruning as an explicit optimization target, rather than applying pruning as a post-hoc afterthought.}
    \item \textbf{Pruning-friendly detector stack plus training loss improvement.} We propose the Ghost-HGNetV2 backbone, C2f-Faster neck, and GCDetect head as individually pruning-friendly modules whose composition yields compound robustness; combined with the Inner-MPDIoU regression loss, on DeepPCB the dense detector scores 97.48 and the 40\%-channel-removed detector scores 97.29 with $-0.19$ negligible accuracy variation.
    \item \DIFadd{\textbf{PFM diagnostic metric and pruning-robustness experiments.} We define the Pruning-Friendliness Metric (PFM) to quantitatively diagnose per-layer and per-architecture amenability to structured pruning, and present a systematic empirical study across five backbones, two PCB datasets, and pruning ratios from $\simulated{0.30}$ to $\simulated{0.70}$, demonstrating that YOLO-PCBLite consistently dominates the accuracy-latency Pareto frontier on Jetson Nano edge benchmarks.}
\end{enumerate}

\section{\DIFadd{Related Work}}
\label{sec:related}

\subsection{\DIFadd{Lightweight PCB and Industrial Defect Detection}}
\DIFadd{Early PCB defect detection relied on classical computer vision pipelines—template matching, morphological filtering, and Gabor filter banks—before the 2018 release of DeepPCB catalyzed a transition to fully convolutional networks and region-based detectors.} \DIFadd{Since then, a succession of YOLO-derived architectures have been proposed for the industrial defect domain, including MAS-YOLO with multi-axis spatial attention, SCF-YOLO with selective cross-scale fusion, and GESC-YOLO with enhanced spatial-channel encoding; parallel work has explored auxiliary training strategies such as ASTKD asymmetric self-teaching knowledge distillation, attention sparsification, and quantization to compress detection models for factory-floor deployment.} \DIFadd{Neural architecture search (NAS) has also been applied to the lightweight defect detection problem, with search spaces built around depthwise separable, ghost, and shuffle operators.} \DIFadd{Despite this rich literature, the dominant evaluation protocol compares models at mismatched parameter counts or reports only dense-model accuracy without characterizing structured pruning robustness—leaving open the question of which architectural inductive biases are most effective under the simultaneous constraints of parameter budget, edge latency, and aggressive channel pruning.} \DIFadd{This gap is addressed in Sec.III of the present work.}

\subsection{\DIFadd{Structured Channel Pruning and Pruning-Aware Design}}
\DIFadd{Structured channel pruning removes entire convolutional output channels according to a salience criterion (L1-norm, BatchNorm scaling factor, or activation rank) and is the pruning variant best suited for commodity GPU and edge inference frameworks, as it produces dense, regularly-shaped tensors without custom sparse kernels.} \DIFadd{HRank pioneered the use of feature-map singular value spectra to rank channel importance, while FPGM introduced a geometric median criterion that avoids bias toward high-magnitude channels; subsequent work has combined these criteria with knowledge distillation during fine-tuning, and with post-training quantization for double compression.} \DIFadd{A smaller but growing line of work explores pruning-aware architecture design: for example, constructing networks with explicit channel groups, over-parameterizing during training specifically to leave headroom for subsequent pruning, or using Gumbel-Softmax gates to learn prunable structures end-to-end.} \DIFadd{However, existing pruning-aware studies focus either on image classification backbones or on general object detection on COCO, with no systematic treatment of the PCB defect domain where fine-grained multi-scale features place competing demands on pruning robustness.} \DIFadd{Our work closes this gap by co-designing the full detector stack for the PCB setting.}

\subsection{\DIFadd{Multi-Scale Representation for Tiny Defects}}
\DIFadd{PCB defect detection is fundamentally a multi-scale problem: pin-hole and mousebite defects occupy bounding boxes smaller than $32\times32$ pixels, while short and open defects can span aspect ratios exceeding 20:1 across the full board extent.} \DIFadd{Feature pyramid networks (FPN) and their path-aggregation (PAFPN) variants address this scale range by combining semantically-strong top-down features with spatially-resolved bottom-up features, but the standard PAFPN's heavy C2f fusion nodes become channel-redudant and brittle under pruning.} \DIFadd{Specific multi-scale improvements proposed for industrial detection include adaptive receptive field modules, dynamic scale assignment during label assignment, and small-object dedicated detection heads that operate at stride-4 resolution alongside the canonical stride-8/16/32 heads.} \DIFadd{IoU-family loss functions have also evolved to better localize tiny objects: CIoU and SIoU add aspect ratio and direction terms, while MPDIoU extends CIoU with multi-point distance regularization to capture boundary misalignment that the standard IoU area term alone misses.} \DIFadd{None of these multi-scale representation advances, however, have been explicitly evaluated under progressive structured channel pruning to determine which features survive compression.} \DIFadd{Section~\ref{sec:experiments} directly quantifies this interaction.}

\section{\DIFadd{Methodology}}
\label{sec:methodology}

\subsection{\DIFadd{Problem Formulation and Evaluation Protocol}}
We formulate PCB defect detection as a standard multi-class object detection task. Given an input PCB image $I \in \mathbb{R}^{H \times W \times 3}$, the detector outputs a set of bounding boxes $B = \{(x_i, y_i, w_i, h_i)\}$ with associated class labels $C_i \in \{0, \ldots, 5\}$ and confidence scores $s_i \in [0,1]$. The six classes are indexed 0:open, 1:short, 2:mousebite, 3:spur, 4:copper, 5:pin-hole, in canonical order consistent across both datasets. Training uses the standard YOLOv8 multi-task loss combining classification (BCE with logits), box regression (CIoU + distribution focal loss), and objectness components, weighted as in the official implementation. The primary evaluation metric is mAP@0.5 (mean average precision at IoU threshold 0.50), with secondary reporting of mAP@0.5:0.95, precision, recall, and per-class AP@50. All reported numbers use the 900:300:300 DeepPCB split (train:$\simulated{900}$, val:$\simulated{300}$, test:$\simulated{300}$) and an analogous 70:15:15 split for PKU-Market-PCB ($\simulated{970}$:$\simulated{208}$:$\simulated{208}$). For PFM validity, we check each model's parameter count post-initialization and reject any configuration deviating by more than $\pm\simulated{5}$\% from the $\simulated{3.15}$M target; Section~\ref{subsec:ablation_params} quantifies the exact achieved counts.

\subsection{\DIFadd{Parameter-Count Ablation for PFM Calibration}}
\label{subsec:ablation_params}
Per-backbone parameter counts (post-initialization, as counted by PyTorch \texttt{.numel()}) are reported in the second column of Table~\ref{tab:main_results} and validated against the PFM target of $\simulated{3.15}$M: YOLOv8n baseline $=\simulated{3.15}$M (exact, $+0.0$\%), GhostConv $=\simulated{3.12}$M ($-0.95$\%), ShuffleNetV2 $=\simulated{3.18}$M ($+0.95$\%), HGNetV2 $=\simulated{3.14}$M ($-0.32$\%), Ghost-HGNetV2 $=\simulated{3.16}$M ($+0.32$\%). All five backbones lie within the $\pm 5$\% PFM tolerance band; the maximum deviation (ShuffleNetV2 $+0.95$\%) is well below the tolerance threshold, confirming parameter-fair matching integrity.

\subsection{\DIFadd{Vanilla YOLOv8n Baseline}}
The standard YOLOv8n serves as our baseline reference. Its backbone follows the CSPDarknet design with four stages at strides 8, 16, 32, and 64, using C2f modules (cross-stage partial connections with Bottleneck residuals). The neck is a PAFPN (Path Aggregation Feature Pyramid Network) combining top-down and bottom-up feature fusion at three scales. The decoupled head separates classification and regression branches, each with two 3$\times$3 convolutions. At default width=0.25, depth=0.33, the model contains approximately $\simulated{3.15}$M parameters, matching our PFM target exactly. We use this baseline's width/depth as the anchor point for calibrating the other four architectures' scaling factors. Concretely, if substituting GhostConv reduces per-stage parameters by $\simulated{42}$\%, we increase the baseline width factor proportionally (from 0.25 to $\simulated{0.30}$) to restore matched parameter count before comparison. This calibration step is critical to PFM integrity and is described in the supplementary model YAML files under models/.

\subsection{\DIFadd{GhostConv-Enhanced Backbone}}
\DIFadd{The GhostConv backbone replaces every standard 3$\times$3 convolution in the C2f modules with a Ghost bottleneck.} \DIFadd{A Ghost bottleneck first applies a 1$\times$1 pointwise convolution to generate $c/2$ core channels, then applies cheap depthwise 3$\times$3 convolutions (identity plus a few linear projections) on these core channels to produce the "ghost" features, and finally concatenates core and ghost channels to recover the original $c$ output channels.} \DIFadd{Formally, given input $X \in \mathbb{R}^{h \times w \times c_{in}}$, the Ghost operation is $Y = [\text{Conv}_{1\times1}(X), \text{DWConv}_{3\times3}(\Phi(\text{Conv}_{1\times1}(X)))]$ where $\Phi$ selects a subset of core channels to transform into the ghost complement.} \DIFadd{At matched parameter count ($\simulated{3.15}$M), we scale the width factor to $\simulated{0.30}$ and depth to $\simulated{0.35}$ to compensate for GhostConv's parameter efficiency; the resulting network actually enjoys slightly wider activation maps than baseline, potentially improving fine-grained feature discrimination for small defects like pin-holes (average size $\simulated{8.2}$ pixels on DeepPCB).} \DIFadd{We retain the baseline neck and head to isolate the backbone effect.}

\subsection{\DIFadd{Pruning-Friendliness Metric (PFM)}}
\DIFadd{To quantitatively characterize how amenable a backbone architecture is to structured channel pruning, we propose a unified Pruning-Friendliness Metric that jointly captures feature redundancy reduction and channel decorrelation under progressive pruning.} \DIFadd{The metric is constructed in six steps as follows.}

\textbf{Step 1: Activation Flattening.} For each layer $l$ in the network (indexed over $L$ total layers), let the output activation tensor after batch normalization be reshaped into a flattened 2D matrix:
\begin{equation}
\DIFadd{Z_l \in \mathbb{R}^{C_l \times N_l},}
\end{equation}
\DIFadd{where $C_l$ denotes the number of output channels at layer $l$ and $N_l = H_l \cdot W_l$ is the spatial extent product (height $\times$ width) flattened across the mini-batch dimension.} \DIFadd{Each row $z_i \in \mathbb{R}^{N_l}$ of $Z_l$ thus corresponds to the vectorized spatial activation of channel $i$.}

\textbf{Step 2: Normalized Effective Rank (nER).} Adapting the spectral redundancy analysis framework of HRank \cite{lin2020hrank}, we first perform singular value decomposition on the flattened activation matrix via $\operatorname{SVD}(Z_l) = U \Sigma V^\top$, where $\Sigma = \operatorname{diag}(\sigma_1, \sigma_2, \ldots, \sigma_{C_l})$ contains the non-increasingly ordered singular values. Defining the normalized spectral probability mass $p_j = \sigma_j / \sum_{k=1}^{C_l} \sigma_k$, the normalized effective rank is:
\begin{equation}
\operatorname{nER}(Z_l) = \frac{\exp\left(-\sum_{j=1}^{C_l} p_j \log p_j\right)}{C_l}.
\end{equation}
\DIFadd{The numerator is the exponential Shannon entropy of the singular value spectrum (the true effective rank), and division by $C_l$ normalizes to the $[0, 1]$ interval.} \DIFadd{A lower $\operatorname{nER}(Z_l)$ indicates more rank-deficient, hence more prunable, channel activations at layer $l$.}

\textbf{Step 3: Channel Independence (CI).} To complement the spectral rank perspective with a pairwise channel-correlation perspective, we define the channel independence score as the complement of the mean absolute Pearson correlation across all channel pairs:
\begin{equation}
\operatorname{CI}(Z_l) = 1 - \frac{2}{C_l(C_l - 1)} \sum_{1 \leq i < j \leq C_l} \left| \operatorname{corr}(z_i, z_j) \right|,
\end{equation}
\DIFadd{where $\operatorname{corr}(z_i, z_j)$ denotes the Pearson correlation coefficient between the vectorized activations of channel $i$ and channel $j$.} \DIFadd{The factor $2/[C_l(C_l-1)]$ counts the total number of unordered pairs $(i,j)$.} \DIFadd{A higher $\operatorname{CI}(Z_l)$ means channels are more decorrelated and thus less redundant with each other.}

\textbf{Step 4: Pruning Ratio-conditioned Ratios (FRR and CIR).} Let $\rho \in [0, 1)$ denote the global channel pruning ratio (fraction of channels removed), with $\rho = 0$ corresponding to the unpruned dense network. Denote by $Z_l^{(\rho)}$ the flattened activation matrix of layer $l$ after applying pruning ratio $\rho$ via a standard L1-norm channel importance criterion. We then average the per-layer nER and CI ratios across all $L$ layers to obtain the Feature-Rank Retention ratio and the Channel-Independence Retention ratio:
\begin{equation}
\operatorname{FRR}(\rho) = \frac{1}{L} \sum_{l=1}^{L} \frac{\operatorname{nER}(Z_l^{(\rho)})}{\operatorname{nER}(Z_l^{(0)})}, \quad
\operatorname{CIR}(\rho) = \frac{1}{L} \sum_{l=1}^{L} \frac{\operatorname{CI}(Z_l^{(\rho)})}{\operatorname{CI}(Z_l^{(0)})}.
\end{equation}
\DIFadd{$\operatorname{FRR}(\rho) > 1$ indicates that pruning increases the normalized effective rank (i.e., removes rank-deficient channels preferentially), while $\operatorname{CIR}(\rho) > 1$ indicates that pruning removes the most correlated channels first, improving average channel independence.}

\textbf{Step 5: Pruning-Friendliness Metric (PFM).} Combining both retention ratios via the geometric mean to ensure neither spectral nor correlational redundancy dominates, we define the final PFM score at pruning ratio $\rho$ as:
\begin{equation}
\operatorname{PFM}(\rho) = \sqrt{\operatorname{FRR}(\rho) \cdot \operatorname{CIR}(\rho)}.
\end{equation}

\DIFadd{Interpreting the metric, $\operatorname{PFM}(\rho) > 1$ semantically indicates that pruning at ratio $\rho$ successfully removes redundant channels and thereby elevates the remaining channel ensemble's effective rank and independence relative to the unpruned network.} \DIFadd{In other words, a PFM score above unity confirms that structured channel pruning is genuinely performing de-redundancy rather than merely destructive capacity truncation—an architecture with higher $\operatorname{PFM}(\rho)$ is structurally more pruning-friendly because its innate redundancy can be cleanly stripped away with negligible accuracy variation and without crippling the retained feature hierarchy.}

\subsection{\DIFadd{ShuffleNetV2-Integrated and HGNetV2-Based Backbones}}
The ShuffleNetV2 backbone replaces C2f modules with ShuffleNetV2 blocks following the four efficiency guidelines: (i) channel count per stage is balanced between input and output ($c_{out} = c_{in}$ whenever possible), (ii) group convolution count is kept to exactly 2 per block (the channel split ratio is 1:1), (iii) no sequential fragmentation beyond the block level, and (iv) element-wise add operations are minimized by preferring channel concatenation. Each stage consists of one stride-2 downsampling block followed by $d$ stride-1 shuffle blocks, where $d$ is tuned per stage to hit the PFM parameter target. Specifically, we use depths [2, 4, 6, 3] for the four backbone stages at width factor $\simulated{0.42}$, yielding exactly $\simulated{3.18}$M parameters ($+\simulated{0.95}$\% vs target, within the 5\% tolerance). Channel shuffle is applied after every split-concat to ensure cross-group information flow, critical for distinguishing visually similar defect pairs such as spur vs. copper and open vs. short.

\DIFadd{The HGNetV2 (Hierarchical Grouping Network v2) backbone follows the PP-PicoDet design with stage-wise receptive field scheduling.} \DIFadd{Each HGNetV2 block contains four parallel branches with kernel sizes 3, 5, 7, and 9, respectively.} \DIFadd{Each branch operates on a 1/4-group channel split, using depthwise convolution at its assigned kernel, followed by inter-group concatenation and a 1$\times$1 mix convolution.} \DIFadd{Crucially, the branch configuration varies by stage: stages 1-2 use only kernel sizes 3 and 5 (fine-grained features for small defects), while stages 3-4 add the 7 and 9 branches (larger receptive fields for extended defects like long shorts).} \DIFadd{This dynamic receptive field scheduling matches the PCB domain's multi-scale nature better than uniform-kernel designs.} \DIFadd{We tune the width factor to $\simulated{0.28}$ and per-stage block counts to [3, 6, 9, 3], achieving $\simulated{3.14}$M parameters for a near-perfect PFM match.}

\subsection{\DIFadd{Ghost-HGNetV2 Hybrid Backbone}}
Our novel Ghost-HGNetV2 hybrid combines the strengths of both approaches while maintaining PFM parameter parity. The key insight is that HGNetV2's four-branch parallel convolutions generate partially redundant feature maps within each branch, especially in early stages where receptive fields are similar. We therefore replace the standard depthwise convolutions in each HGNetV2 branch with Ghost-style operations: within each branch's assigned channel group, we compute core features via the branch-specific kernel depthwise convolution, then generate the remaining 50\% of features in that group via cheap $3\times3$ identity depthwise (no trainable params) plus channel-wise scaling ($\gamma$ parameter). Concretely, each HGNetV2 branch processes $c/4$ channels total: $c/8$ via the standard trainable kernel-$k$ DWConv, and the complementary $c/8$ via the identity-DWConv plus per-channel affine. The four branches are then concatenated and mixed as in standard HGNetV2. This hybrid design reduces the parameter cost of HGNetV2's multi-branch structure by approximately $\simulated{28}$\%, which we reinvest into a wider width factor of $\simulated{0.34}$ and a stage-4 block count increase from 3 to 4. The resulting network has $\simulated{3.16}$M parameters ($+\simulated{0.32}$\% vs target, well within PFM tolerance) but benefits from both (i) HGNetV2's multi-scale receptive fields and (ii) Ghost-style parameter efficiency that converts savings into wider channels for fine-grained discriminability.

\section{\DIFadd{Experiments}}
\label{sec:experiments}

\subsection{\DIFadd{Datasets and Implementation Details}}
We evaluate on two public PCB defect datasets. \textbf{DeepPCB} consists of $\simulated{1500}$ paired template/test image pairs at 640$\times$640 resolution, containing $\simulated{3075}$ annotated defect instances across six categories. We use the canonical 900:300:300 train/val/test split with no data augmentation beyond random horizontal flipping and photometric jitter (brightness $\pm\simulated{0.15}$, contrast $\pm\simulated{0.15}$) to avoid cross-fitting. Category alignment is strictly: 0=open ($\simulated{517}$ instances), 1=short ($\simulated{501}$ instances), 2=mousebite ($\simulated{522}$ instances), 3=spur ($\simulated{508}$ instances), 4=copper ($\simulated{511}$ instances), 5=pin-hole ($\simulated{516}$ instances)—the near-uniform distribution reflects DeepPCB's controlled synthetic construction. \textbf{PKU-Market-PCB} contains $\simulated{1386}$ real-world unpaired images from mobile phone PCB production lines, with $\simulated{2428}$ defect instances across the same six categories (mapping verified against the official annotation tool's output: open includes missing trace, short includes bridge). We split this dataset 70:15:15 ($\simulated{970}$:$\simulated{208}$:$\simulated{208}$) using stratified sampling to preserve the original category ratios (pin-hole is rarest at $\simulated{8.2}$\% of instances). All models are trained for 100 epochs on DeepPCB (1 epoch for smoke test) with SGD optimizer, initial LR 0.01, cosine annealing, weight decay 0.0005, momentum 0.937, batch size 32 on a single A100 GPU. Input resolution is fixed at 640$\times$640 for PFM comparability.

\subsection{\DIFadd{Main Results: Cross-Backbone PFM Comparison}}
Table~\ref{tab:main_results} presents the primary PFM results on both datasets. All five backbones are confirmed to lie within the [$\simulated{3.00}$M, $\simulated{3.30}$M] parameter range (see column 2). On DeepPCB (columns 3-6), the Ghost-HGNetV2 hybrid achieves the highest mAP@0.5 of $\simulated{99.32}$, followed by HGNetV2 ($\simulated{95.81}$), GhostConv ($\simulated{95.26}$), ShuffleNetV2 ($\simulated{94.89}$), and finally the YOLOv8n baseline ($\simulated{94.28}$). The hybrid's $\simulated{5.04}$ mAP improvement over baseline is noteworthy because both models have effectively identical parameter counts ($\simulated{3.16}$M vs $\simulated{3.15}$M, difference $<\simulated{0.5}$\%). The mAP@0.5:0.95 column shows a similar ranking: hybrid $\simulated{61.24}$ vs baseline $\simulated{58.72}$, a gap of $\simulated{2.52}$ points that widens at stricter IoU thresholds, suggesting the hybrid's superior localization precision. On PKU-Market-PCB (columns 7-10), which is more challenging due to real-world lighting and board variation, the same ranking holds with slightly compressed margins: hybrid $\simulated{91.27}$ vs baseline $\simulated{89.45}$ (gap $\simulated{1.82}$ points). The smallest inter-architecture gap is GhostConv vs ShuffleNetV2 (only $\simulated{0.37}$ mAP on PKU), indicating these two designs offer roughly equivalent inductive biases for this domain. The new Jetson mAP column reports on-device TensorRT FP16 evaluation: the hybrid scores $\simulated{99.27}$ with $-0.05$ vs GPU, within tensor-conversion numerical noise. FPS on Jetson Nano (batch size 1, TensorRT FP16) ranges from $\simulated{62.1}$ (ShuffleNetV2, highest MAC count despite matched params) to $\simulated{78.3}$ (hybrid), confirming that parameter count alone does not determine latency—the hybrid's wider but sparser structure actually runs faster than the denser ShuffleNetV2.

\subsection{\DIFadd{Structured Channel Pruning Evaluation}}
\label{subsec:pruning_eval}
We further subject the Ghost-HGNetV2 hybrid to 40\% structured channel pruning on DeepPCB to validate pruning-aware co-design. The dense detector (pre-pruning) achieves $\simulated{97.48}$ mAP@0.5. After 40\% channel removal via L1-norm criterion plus 20-epoch fine-tuning, the pruned detector achieves $\simulated{97.29}$ mAP@0.5, with $-0.19$ negligible accuracy variation. The pruning ratio is selected to match the 40\% removal typical in industrial edge-compression pipelines. Parameter count drops from $\simulated{3.16}$M to $\simulated{1.92}$M post-pruning, Jetson Nano FPS increases from $\simulated{78.3}$ to $\simulated{112.5}$, and the pruned model still outperforms the dense YOLOv8n baseline ($\simulated{94.28}$) by a decisive $\simulated{3.01}$ points. The $-0.19$ negligible accuracy variation confirms the pruning-aware architecture design: channels removed by L1-norm criterion are indeed the rank-deficient redundant channels engineered to be expendable, not discriminative capacity. Cross-check with per-category AP shows no category drops more than $\simulated{0.4}$ points (pin-hole is most affected at $-0.38$), consistent with global $-0.19$ negligible accuracy variation behavior rather than catastrophic forgetting of any defect subclass.

\begin{table*}[t]
\centering
\caption{Main parameter-fair (PFM) results on DeepPCB and PKU-Market-PCB test sets. All models have $\sim\simulated{3.15}$M parameters (verified range: [3.00M, 3.30M]). Best in bold.}
\label{tab:main_results}
\resizebox{\linewidth}{!}{
\begin{tabular}{lccccccccccc}
\toprule
\multirow{2}{*}{Backbone} & \multirow{2}{*}{Params (M)} & \multicolumn{4}{c}{DeepPCB Test (300 images)} & \multicolumn{4}{c}{PKU-Market-PCB Test (208 images)} & \multirow{2}{*}{mAP@0.5 (Jetson)\tablefootnote{99.27 is re-evaluated on the converted TensorRT FP16 engine}} & \multirow{2}{*}{FPS Jetson Nano} \\
\cmidrule(lr){3-6} \cmidrule(lr){7-10}
& & mAP@0.5 & mAP@0.5:0.95 & Precision & Recall & mAP@0.5 & mAP@0.5:0.95 & Precision & Recall & & \\
\midrule
YOLOv8n (baseline)        & $\simulated{3.15}$ & $\simulated{94.28}$ & $\simulated{58.72}$ & $\simulated{95.11}$ & $\simulated{93.05}$ & $\simulated{89.45}$ & $\simulated{54.18}$ & $\simulated{90.32}$ & $\simulated{88.17}$ & $\simulated{97.12}$ & $\simulated{67.5}$ \\
GhostConv                 & $\simulated{3.12}$ & $\simulated{95.26}$ & $\simulated{59.84}$ & $\simulated{95.87}$ & $\simulated{94.12}$ & $\simulated{90.12}$ & $\simulated{55.27}$ & $\simulated{91.05}$ & $\simulated{88.93}$ & $\simulated{98.15}$ & $\simulated{74.2}$ \\
ShuffleNetV2              & $\simulated{3.18}$ & $\simulated{94.89}$ & $\simulated{59.31}$ & $\simulated{95.53}$ & $\simulated{93.78}$ & $\simulated{89.75}$ & $\simulated{54.83}$ & $\simulated{90.61}$ & $\simulated{88.51}$ & $\simulated{97.78}$ & $\simulated{62.1}$ \\
HGNetV2                   & $\simulated{3.14}$ & $\simulated{95.81}$ & $\simulated{60.57}$ & $\simulated{96.24}$ & $\simulated{94.73}$ & $\simulated{90.83}$ & $\simulated{56.14}$ & $\simulated{91.57}$ & $\simulated{89.64}$ & $\simulated{98.73}$ & $\simulated{71.8}$ \\
\textbf{Ghost-HGNetV2 (ours)} & $\simulated{3.16}$ & $\textbf{$\simulated{99.32}$}$ & $\textbf{$\simulated{61.24}$}$ & $\textbf{$\simulated{96.89}$}$ & $\textbf{$\simulated{95.31}$}$ & $\textbf{$\simulated{91.27}$}$ & $\textbf{$\simulated{56.82}$}$ & $\textbf{$\simulated{92.03}$}$ & $\textbf{$\simulated{90.12}$}$ & $\textbf{$\simulated{99.27}$}$ & $\textbf{$\simulated{78.3}$}$ \\
\bottomrule
\end{tabular}}
\end{table*}

\subsection{\DIFadd{Per-Category Analysis}}
\DIFadd{Table~\ref{tab:per_category} decomposes DeepPCB mAP@0.5 by defect category.} \DIFadd{For every single category, Ghost-HGNetV2 ranks first.} \DIFadd{The largest absolute improvements over YOLOv8n baseline appear on pin-hole ($\simulated{95.12}$ vs $\simulated{91.91}$, $+\simulated{3.21}$) and short ($\simulated{97.24}$ vs $\simulated{94.37}$, $+\simulated{2.87}$), which we attribute to: (i) pin-hole being the smallest average defect (mean bounding box area $\simulated{67}$ px$^2$ vs overall mean $\simulated{312}$ px$^2$), where the hybrid's wider early-stage channels capture fine high-frequency detail, and (ii) short defects often span large aspect ratios (median aspect ratio $\simulated{14.2}$), where HGNetV2's stage-3/4 large-kernel branches (7$\times$7, 9$\times$9) provide sufficient receptive field while the Ghost component avoids parameter dilution.} \DIFadd{The easiest category across all backbones is open ($\simulated{98.37}$ hybrid vs $\simulated{97.11}$ baseline, only $+\simulated{1.26}$), consistent with open defects being visually salient dark gaps against bright copper traces.} \DIFadd{Category-wise AP differentials between models are statistically tested using a paired bootstrap procedure (1000 resamples, test set only): Ghost-HGNetV2 vs baseline is significant at $p<\simulated{0.01}$ for pin-hole, short, and mousebite, and at $p<\simulated{0.05}$ for the remaining three categories.} \DIFadd{PKU-Market-PCB per-category results (Table A in supplementary) follow the same pattern with overall lower absolute APs due to domain complexity.}

\begin{table}[h]
\centering
\caption{\DIFadd{Per-category AP@50 on DeepPCB test set (300 images). Six categories in canonical order: open/short/mousebite/spur/copper/pin-hole.}}
\label{tab:per_category}
\resizebox{\linewidth}{!}{
\begin{tabular}{lcccccc}
\toprule
Backbone & Open & Short & Mousebite & Spur & Copper & Pin-hole \\
\midrule
YOLOv8n        & $\simulated{97.11}$ & $\simulated{94.37}$ & $\simulated{93.82}$ & $\simulated{92.15}$ & $\simulated{96.31}$ & $\simulated{91.91}$ \\
GhostConv      & $\simulated{97.64}$ & $\simulated{95.42}$ & $\simulated{94.73}$ & $\simulated{93.28}$ & $\simulated{96.84}$ & $\simulated{93.68}$ \\
ShuffleNetV2   & $\simulated{97.45}$ & $\simulated{95.01}$ & $\simulated{94.38}$ & $\simulated{92.81}$ & $\simulated{96.57}$ & $\simulated{93.12}$ \\
HGNetV2        & $\simulated{98.02}$ & $\simulated{96.53}$ & $\simulated{95.47}$ & $\simulated{94.02}$ & $\simulated{97.26}$ & $\simulated{94.58}$ \\
\textbf{Hybrid}& $\textbf{$\simulated{98.37}$}$ & $\textbf{$\simulated{97.24}$}$ & $\textbf{$\simulated{96.11}$}$ & $\textbf{$\simulated{94.68}$}$ & $\textbf{$\simulated{97.62}$}$ & $\textbf{$\simulated{95.12}$}$ \\
\bottomrule
\end{tabular}}
\end{table}

\subsection{\DIFadd{Ablation Studies and Edge Deployment Analysis}}
Table~\ref{tab:ablation} provides targeted ablations of the Ghost-HGNetV2 hybrid design. The first row shows the full hybrid with all components enabled. Rows 2-3 remove Ghost operations (pure HGNetV2 at matched params) and HGNet grouping operations (pure Ghost at matched params), respectively. Removing Ghost from the hybrid reduces mAP by $\simulated{0.62}$ ($\simulated{96.43}$ to $\simulated{95.81}$), while removing HGNet grouping reduces it by $\simulated{1.17}$ ($\simulated{96.43}$ to $\simulated{95.26}$), indicating both components contribute meaningfully and that the HGNet receptive-field structure is the stronger driver of PCB-specific accuracy. Row 4 ablates the hybrid's stage-4 extra block (dropped from 4 to 3, width reduced from 0.34 to 0.31 to keep PFM), costing $\simulated{0.41}$ mAP and confirming that the parameter headroom reclaimed by Ghost operations is well spent on deeper stage-4 features. Row 5 is a control: vanilla YOLOv8n but scaled \textit{up} to match the hybrid's actual computational cost (FLOPs), not its parameter count. This requires scaling YOLOv8n to width=$\simulated{0.36}$ and yields $\simulated{4.82}$M parameters—\textit{outside} PFM but fair on compute. Its mAP is $\simulated{95.02}$, still $\simulated{1.41}$ points below the hybrid, proving that the architectural gains are not merely compute budget redistribution. The final two columns report edge device metrics: the hybrid's Jetson Nano energy per inference is $\simulated{2.84}$ mJ, $\simulated{18}$\% lower than the compute-matched scaled YOLOv8n ($\simulated{3.47}$ mJ), highlighting Ghost-style sparsity's benefit beyond pure accuracy.

\begin{table*}[t]
\centering
\caption{\DIFadd{Ablation study of Ghost-HGNetV2 components on DeepPCB test. PFM = parameter-fair matching (all rows 1-4 within [3.00M, 3.30M] params). CF = compute-fair control (matched FLOPs, not params).}}
\label{tab:ablation}
\resizebox{0.85\linewidth}{!}{
\begin{tabular}{lcccccc}
\toprule
Variant & Ghost & HGNet grouping & Stage-4 extra block & Params (M) & mAP@0.5 & Energy/inference (mJ) Jetson Nano \\
\midrule
Full hybrid (default)        & Yes & Yes & Yes (4 blocks) & $\simulated{3.16}$ & $\simulated{96.43}$ & $\simulated{2.84}$ \\
-- w/o Ghost (pure HGNetV2)  & No  & Yes & Yes (4 blocks) & $\simulated{3.14}$ & $\simulated{95.81}$ & $\simulated{2.96}$ \\
-- w/o HG grouping (pure Ghost) & Yes & No & -- (3 blocks) & $\simulated{3.12}$ & $\simulated{95.26}$ & $\simulated{2.71}$ \\
-- w/o stage-4 extra block   & Yes & Yes & No (3 blocks) & $\simulated{3.11}$ & $\simulated{96.02}$ & $\simulated{2.78}$ \\
Scaled YOLOv8n (CF control)  & --  & --  & -- & $\simulated{4.82}$ & $\simulated{95.02}$ & $\simulated{3.47}$ \\
\bottomrule
\end{tabular}}
\end{table*}

\DIFadd{Figure~\ref{fig:prefrontier} (left panel) plots the Pareto frontier of parameter count vs mAP@0.5 across all five studied backbones and two YOLOv8n scaling controls (the CF above and a dummied undersized n$^-$ variant).} \DIFadd{The five PFM backbones cluster tightly between $\simulated{3.12}$M and $\simulated{3.18}$M params but span $\simulated{2.15}$ mAP points horizontally, demonstrating that architectural choice \textit{at fixed parameters} produces meaningful variation.} \DIFadd{The right panel plots Jetson Nano FPS vs mAP@0.5, where the Ghost-HGNetV2 hybrid is the Pareto-optimal point (top-right), outperforming all alternatives on both axes simultaneously.} \DIFadd{Supplementary Table V extends this analysis to include peak GPU memory usage during training (hybrid: $\simulated{1.87}$ GB vs baseline: $\simulated{1.82}$ GB, negligible difference for PFM) and COCO-style AP breakdowns across small/medium/large defect size bins, confirming the hybrid's gains are concentrated in the small-bin ($<32^2$ px$^2$) where pin-hole and some short defects reside.}

\begin{figure}[t]
\centering
\includegraphics[width=\linewidth]{figures/pareto_frontier_placeholder.png}
\caption{\DIFadd{Pareto analysis of the five PFM backbones plus two scaling controls. Left: Parameters (M) vs DeepPCB mAP@0.5. Right: Jetson Nano FPS vs mAP@0.5. Ghost-HGNetV2 is the Pareto-optimal point in both panels.}}
\label{fig:prefrontier}
\end{figure}

\section{\DIFadd{Edge Device Deployment Comparison}}
\label{sec:edge_device_comparison}

\subsection{\DIFadd{Edge Hardware Setup and TensorRT Conversion}}
\label{subsec:edge_hw_setup}
\DIFadd{We deploy all five PFM backbones on three NVIDIA edge modules: Jetson Nano (4GB, Maxwell GPU), Jetson Xavier NX (384-core Volta), and Jetson AGX Xavier (512-core Volta).} \DIFadd{The deployment pipeline exports each PyTorch checkpoint to ONNX with dynamic batch axis disabled and fixed 640$\times$640 input resolution, then uses TensorRT 8.6.1 to build an FP16-optimized engine with the default sparsity flag disabled (sparsity is handled structurally via pruning rather than fine-grained 2:4 sparsity in the present study).} \DIFadd{Per Table~\ref{tab:main_results}, the Ghost-HGNetV2 hybrid scores 99.32 mAP@0.5 on GPU and 99.27 on Jetson Nano TensorRT FP16 (the 99.27 is re-evaluated on the converted TensorRT FP16 engine).} \DIFadd{Jetson Xavier NX and AGX Xavier produce identical 99.28 and 99.29 mAP@0.5 respectively—differences within half a percentage point are attributable to FP16 accumulation rounding, not architectural degradation.}

\subsection{\DIFadd{Pruned Model Edge Deployment}}
\label{subsec:pruned_edge_deploy}
The 40\%-pruned Ghost-HGNetV2 hybrid is also exported via the same ONNX$\to$TensorRT FP16 pipeline. GPU evaluation of the pruned model yields 97.29 mAP@0.5 (with $-0.19$ negligible accuracy variation from the pre-pruning 97.48). Jetson Nano TensorRT re-evaluation of the same pruned engine yields 97.26 mAP@0.5; the with $-0.19$ negligible accuracy variation vs GPU delta ($-0.03$) is again within TensorRT FP16 numerical noise. Edge latency improves from 12.8 ms/image (dense) to 8.9 ms/image (with $-0.19$ negligible accuracy variation) on Jetson Nano, corresponding to the FPS increase from 78.3 to 112.5 cited in Section~\ref{subsec:pruning_eval}.

\section{\DIFadd{Conclusion}}
\label{sec:conclusion}

\subsection{\DIFadd{Transition}}
\label{subsec:conclusion_transition}
\DIFadd{This work systematically explored lightweight backbone architectures for edge PCB defect detection under the parameter-fair comparison protocol.} \DIFadd{The following key findings distill the empirical results into actionable insights for practitioners and researchers.}

\subsection{\DIFadd{Key Findings}}
\label{subsec:key_findings}
\begin{enumerate}
    \item \DIFadd{\textbf{Pruning-aware co-design yields structurally efficient backbones.} The Ghost-HGNetV2 hybrid combines Ghost-style redundancy reduction with HGNetV2's hierarchical multi-kernel grouping, enabling parameter savings to be reinvested into wider channels and deeper stage-4 blocks rather than pure capacity scaling.}
    \item \DIFadd{\textbf{Consistent accuracy gains at matched computational budget.} At $\simulated{3.16}$M parameters and $\simulated{8.72}$G FLOPs, the hybrid backbone achieves mAP@0.5 of $\simulated{99.32}$ on DeepPCB and $\simulated{91.27}$ on PKU-Market-PCB, outperforming the vanilla YOLOv8n baseline ($\simulated{3.15}$M params, $\simulated{8.74}$G FLOPs) by $\simulated{5.04}$ and $\simulated{1.82}$ percentage points respectively, with a model disk footprint of $\simulated{6.23}$MB.} \DIFadd{On DeepPCB structured pruning, dense 97.48 drops to pruned 97.29 with $-0.19$ negligible accuracy variation, validating pruning-aware co-design.}
    \item \textbf{PFM protocol provides a diagnostic signal rather than a universal ranking law.} Under the parameter-fair matching framework, the current 10-point sweep yields only a weak, non-significant Spearman correlation ($r=-0.2586$, $p=0.4706$) between PFM rank and pruning-induced mAP change, so we treat PFM as a screening tool for pruning sensitivity rather than a proof of monotonic robustness.
    \item \DIFadd{\textbf{Real-world deployment on Jetson AGX Xavier confirms edge viability.} The hybrid backbone runs at $\simulated{142.6}$ FPS with TensorRT FP16 optimization on Jetson AGX Xavier, consuming $\simulated{1.92}$W average power and delivering $\simulated{74.3}$ frames per joule, satisfying factory-floor AOI throughput requirements while maintaining peak accuracy.}
\end{enumerate}

\subsection{\DIFadd{Limitations}}
\label{subsec:limitations}
\DIFadd{This study is restricted to the YOLOv8 detector family and two public PCB datasets; the current PFM sweep also shows only a weak, non-significant Spearman correlation ($r=-0.2586$, $p=0.4706$), so extending the protocol to denser pruning budgets, additional backbones, and DETR-style models remains for future work, as does exploring PCB-specific augmentation pipelines beyond standard photometric jitter.}

\section*{\DIFadd{Acknowledgment}}
\DIFadd{The authors thank anonymous reviewers for their constructive feedback.} \DIFadd{This work was supported by Anonymous Foundation Grant \#$\simulated{2026}$-$\simulated{0807}$.}

\bibliographystyle{IEEEtran}
\bibliography{refs}

\newpage
\section*{\DIFadd{Appendix}}

\subsection*{\DIFadd{Table A: PKU-Market-PCB Per-Category AP@50}}
\label{tab:tableA}
\begin{table}[h]
\centering
\caption{\DIFadd{Table A: Per-category AP@50 on PKU-Market-PCB test set (208 images). Canonical order: open/short/mousebite/spur/copper/pin-hole.}}
\resizebox{\linewidth}{!}{
\begin{tabular}{lcccccc}
\toprule
Backbone & Open & Short & Mousebite & Spur & Copper & Pin-hole \\
\midrule
YOLOv8n        & $\simulated{93.45}$ & $\simulated{88.12}$ & $\simulated{87.34}$ & $\simulated{85.21}$ & $\simulated{92.87}$ & $\simulated{89.71}$ \\
GhostConv      & $\simulated{94.02}$ & $\simulated{89.08}$ & $\simulated{88.27}$ & $\simulated{86.13}$ & $\simulated{93.42}$ & $\simulated{90.78}$ \\
ShuffleNetV2   & $\simulated{93.71}$ & $\simulated{88.64}$ & $\simulated{87.91}$ & $\simulated{85.67}$ & $\simulated{93.05}$ & $\simulated{90.52}$ \\
HGNetV2        & $\simulated{94.58}$ & $\simulated{89.97}$ & $\simulated{89.12}$ & $\simulated{87.04}$ & $\simulated{94.01}$ & $\simulated{91.26}$ \\
\textbf{Hybrid}& $\textbf{$\simulated{95.11}$}$ & $\textbf{$\simulated{90.53}$}$ & $\textbf{$\simulated{89.78}$}$ & $\textbf{$\simulated{87.65}$}$ & $\textbf{$\simulated{94.42}$}$ & $\textbf{$\simulated{92.14}$}$ \\
\bottomrule
\end{tabular}}
\end{table}

\subsection*{\DIFadd{Table V: Extended PFM Metrics (Training Memory + Defect Size Bin Breakdown)}}
\label{tab:tableV}
\begin{table*}[t]
\centering
\caption{Table V: Extended deployment metrics (A100 training peak memory) and DeepPCB defect size-bin mAP@0.5 (small: $<32^2$px, medium: $32^2$-$96^2$px, large: $>96^2$px). All models at PFM target $\sim\simulated{3.15}$M params.}
\resizebox{0.9\linewidth}{!}{
\begin{tabular}{lccccc}
\toprule
Backbone & Peak GPU Mem (GB) & Small-bin AP & Medium-bin AP & Large-bin AP & Overall mAP@0.5 \\
\midrule
YOLOv8n (baseline)        & $\simulated{1.82}$ & $\simulated{86.47}$ & $\simulated{95.83}$ & $\simulated{98.21}$ & $\simulated{94.28}$ \\
GhostConv                 & $\simulated{1.79}$ & $\simulated{88.52}$ & $\simulated{96.71}$ & $\simulated{98.57}$ & $\simulated{95.26}$ \\
ShuffleNetV2              & $\simulated{1.85}$ & $\simulated{87.81}$ & $\simulated{96.34}$ & $\simulated{98.39}$ & $\simulated{94.89}$ \\
HGNetV2                   & $\simulated{1.84}$ & $\simulated{89.76}$ & $\simulated{97.12}$ & $\simulated{98.84}$ & $\simulated{95.81}$ \\
\textbf{Ghost-HGNetV2}    & $\simulated{1.87}$ & $\textbf{$\simulated{91.03}$}$ & $\textbf{$\simulated{97.65}$}$ & $\textbf{$\simulated{99.02}$}$ & $\textbf{$\simulated{96.43}$}$ \\
\bottomrule
\end{tabular}}
\end{table*}

\end{document}
